% Chương 5 - Phụ trách: Vũ Đức Minh | Báo cáo ~2-3 trang
\section{ĐÁNH GIÁ TỔNG HỢP}\label{chap:danh-gia}
\canviet{Phụ trách: Vũ Đức Minh. Dung lượng dự kiến: 2--3 trang. Mở chương bằng 2--3 câu giới thiệu nội dung chương.}

\subsection{So sánh Random Forest và SVM}\label{sec:rf-vs-svm}
\canviet{Lập bảng ACC của RF và SVM theo từng thí nghiệm (Bảng 4--10 và Phụ lục A của bài báo). RF tốt hơn trong hầu hết trường hợp, ngoại trừ thí nghiệm seed chưa biết; giải thích vì sao FANCI chọn một RF duy nhất.}

\subsection{So sánh FANCI với các hệ thống liên quan}\label{sec:fanci-vs-khac}
\canviet{Dựa trên Bảng~\ref{tab:so-sanh-phuong-phap}. Lưu ý so sánh khó công bằng vì các hệ thống có mục tiêu và dữ liệu khác nhau, dữ liệu và mã nguồn không công khai.}

\subsection{Ưu điểm của FANCI}\label{sec:uu-diem}
\canviet{Chỉ cần một NXD, không theo dõi lưu lượng; độ chính xác cao, FPR thấp; tổng quát hoá tốt giữa các mạng; nhanh; triển khai được dạng dịch vụ; phát hiện được DGA mới trong thực tế.}

\subsection{Hạn chế của FANCI}\label{sec:han-che}
\canviet{DGA dựa trên danh sách từ (Matsnu, GozNym giai đoạn 2) khó phát hiện hơn; phụ thuộc dữ liệu gán nhãn từ DGArchive; kết quả dương tính trong thực tế vẫn cần kiểm tra thủ công; nguy cơ DGA đối kháng; danh sách trắng Alexa đã ngừng cung cấp từ 2022.}

\subsection{Phương pháp nào tối ưu và vì sao}\label{sec:toi-uu}
\canviet{Kết luận theo mục 5 của README: FANCI tối ưu về sự cân bằng giữa độ chính xác, hiệu năng và khả năng triển khai, nhờ đặc trưng chỉ trích xuất từ một tên miền, không cần ngữ cảnh.}

\subsection{Tiểu kết chương}
\canviet{Tóm tắt đánh giá, dẫn sang Chương~\ref{chap:chuyen-doi-so}.}
